%%%%%%%%%%%%%%%%%%%%%%%%%%%%%%%%%%%%%%%%%%%%%%%
% 공통 매크로 정의 (definitions)
%%%%%%%%%%%%%%%%%%%%%%%%%%%%%%%%%%%%%%%%%%%%%%%

% 양자 표기 (quantikz가 \ket 등을 이미 정의하므로 provide+renew로 안전 재정의)
\providecommand{\ket}[1]{}\renewcommand{\ket}[1]{\left|#1\right\rangle}
\providecommand{\bra}[1]{}\renewcommand{\bra}[1]{\left\langle#1\right|}
\providecommand{\braket}[2]{}\renewcommand{\braket}[2]{\left\langle#1\middle|#2\right\rangle}
\providecommand{\expect}[1]{}\renewcommand{\expect}[1]{\left\langle#1\right\rangle}
\providecommand{\Tr}{}\renewcommand{\Tr}{\mathrm{Tr}}

% Pauli / 게이트
\newcommand{\PauliX}{X}
\newcommand{\PauliY}{Y}
\newcommand{\PauliZ}{Z}
\newcommand{\Had}{H}

% 자주 쓰는 약어
\newcommand{\nq}{n_q}        % qubit 수
\newcommand{\PQC}{PQC}

% 리뷰 대응 수정 표시 (최종 출판본에서는 아래 줄을 \newcommand{\rev}[1]{#1}로 교체)
% \newcommand{\rev}[1]{\textcolor{blue}{#1}}
\newcommand{\rev}[1]{#1}

% 표/그림 경로
\graphicspath{{figs/}}

% 본문 내 그림/표 인용: [그림 1], [표 1] 형식 (조사가 필요하면 \josaref 사용)
\newcommand{\figref}[1]{[그림~\ref{#1}]}
\newcommand{\tabref}[1]{[표~\ref{#1}]}

% ===== 표/그림 번호 발음 기반 자동 조사 =====
% \josaref{표}{label}{UN}  ->  [표 N] + 은/는 (번호가 바뀌어도 자동 일치)
%   UN: 은/는, I: 이/가, UL: 을/를, GW: 과/와
% 끝자리 발음에 받침이 있으면(0,1,3,6,7,8) 받침용 조사, 없으면(2,4,5,9) 모음용 조사.
\makeatletter
\newif\if@jong
\def\jUN{UN}\def\jI{I}\def\jUL{UL}\def\jGW{GW}
\newcommand{\josaref}[3]{%
  [#1~\ref{#2}]%
  \@tempcnta=\getrefnumber{#2}\relax
  \@tempcntb=\@tempcnta \divide\@tempcntb by 10 \multiply\@tempcntb by 10
  \advance\@tempcnta by -\@tempcntb % 끝자리
  \@jongfalse
  \ifnum\@tempcnta=0 \@jongtrue\fi
  \ifnum\@tempcnta=1 \@jongtrue\fi
  \ifnum\@tempcnta=3 \@jongtrue\fi
  \ifnum\@tempcnta=6 \@jongtrue\fi
  \ifnum\@tempcnta=7 \@jongtrue\fi
  \ifnum\@tempcnta=8 \@jongtrue\fi
  \def\@jt{#3}%
  \if@jong
    \ifx\@jt\jUN 은\else\ifx\@jt\jI 이\else\ifx\@jt\jUL 을\else\ifx\@jt\jGW 과\fi\fi\fi\fi
  \else
    \ifx\@jt\jUN 는\else\ifx\@jt\jI 가\else\ifx\@jt\jUL 를\else\ifx\@jt\jGW 와\fi\fi\fi\fi
  \fi
}
\makeatother
